\documentclass{ximera}
\title{Derivatives}
\begin{document}
\begin{abstract}
The slope of a tangent line as a limit of secant slopes, in both the $x\to a$ and the $h\to 0$ forms, leading to the definition of the derivative of $f$ at a number $a$. Textbook §2.7.
\end{abstract}
\maketitle

Remember how at the beginning of limits we looked at a tangent line to a curve. We can now use all the technology we've developed to answer this question for any arbitrary curve!

Let's look at that very first exercise. Let $f(x)=x^{2}$, what is the slope of the tangent line at $x=1$?
\begin{image}
\includegraphics[width=0.6\textwidth]{b48a5324-93dc-45ee-b7c2-dc0b5706bda6-048_806_1051_428_535}
\end{image}

Recall that the method we used was to find secants to the curve. So we used equations like:
\[
m_{\text{sec}}=\frac{f(x)-f(1)}{x-1}=\answer{\frac{x^{2}-1}{x-1}}
\]

If we keep moving $x$ closer and closer to $a$ then we will eventually find the tangent line. In other words, the tangent line at point $(a, f(a))$ has as slope:
\[
m=\lim_{x\rightarrow a}\answer{\frac{f(x)-f(a)}{x-a}}
\]
and from there you can figure out the tangent line for any given point!

\begin{exercise}\label{exercise:4-9}
% AUTHOR CHOICE: example chosen during conversion, change freely
Use this method to find the tangent line $y=m x+b$ when $a=2$.
% work: f(x)=x^2, m = lim_{x->2} (x^2-4)/(x-2) = lim (x+2) = 4; line through (2,4): 4 = 4*2 + b so b = -4
\[
m=\answer{4}, \qquad b=\answer{-4}
\]
\end{exercise}

A lot of books (like ours) find the slope in another method. They fix a distance away from the point $a$ and they try and decrease this distance. What does this mean?
\begin{image}
\includegraphics[width=0.6\textwidth]{b48a5324-93dc-45ee-b7c2-dc0b5706bda6-049_810_1048_1465_537}
\end{image}

Start with some point $a$ where you want to find the tangent line. You let $h$ be some distance away from the point $a$. Then you try and decrease the distance and solve the following limit:
\[
m=\lim_{h\rightarrow 0}\answer{\frac{f(a+h)-f(a)}{h}}
\]

Sometimes, this method is easier to solve than the previous version.

\begin{example}\label{example:4-10}
Let $f(x)=1 / x^{2}$ and find the slope of the tangent line at $\left(2, \frac{1}{4}\right)$.
% work: (1/(2+h)^2 - 1/4)/h = (4-(2+h)^2)/(4h(2+h)^2) = (-4h-h^2)/(4h(2+h)^2) = -(4+h)/(4(2+h)^2) -> -4/16 = -1/4
\[
m=\lim_{h\rightarrow 0}\frac{f(2+h)-f(2)}{h}=\lim_{h\rightarrow 0}\frac{\frac{1}{(2+h)^{2}}-\frac{1}{4}}{h}=\answer{-\frac{1}{4}}
\]
\end{example}

This slope line has a special name and it is (kinda) the whole point of this course.

\begin{definition}\label{definition:4-11}
The \dfn{derivative of a function $f$ at a number $a$} is
\[
f^{\prime}(a)=\lim_{h\rightarrow 0}\answer{\frac{f(a+h)-f(a)}{h}}
\]
\end{definition}

\end{document}
